% "{'chapitre':'slci_multiphy','classe':('PSI'),'type':('application'),'titre':'Mise à l\\'eau d\\'un robot sous-marin','source':'Concours Centrale - MP 2019','comp':['B2-02',],'corrige':True}"

%\setchapterimage{fig_00.jpg}
\chapter*{Application \arabic{cptApplication} \\ 
Etude d’une tuyère à ouverture variable pour banc d’essais de turboréacteur -- \ifprof Corrigé \else Sujet \fi}
\addcontentsline{toc}{section}{Application \arabic{cptApplication} : Etude d’une tuyère à ouverture variable pour banc d’essais de turboréacteur -- \ifprof Corrigé \else Sujet \fi}

\iflivret \stepcounter{cptApplication} \else
\ifprof  \stepcounter{cptApplication} \else \fi
\fi

\setcounter{question}{0}
\marginnote{Concours Banque PT.}





\subsection*{Question de cours}
\question{Donner la strucutre de la chaîne de puissance. On précisera en particulier le nom de chacune des fonctions.}

\question{Donner 3 exemples de composants permettant de réaliser chacune des fonctions.}

\subsection*{Etude d’une tuyère à ouverture variable}
\marginnote{D'après ressources de Florestan Mathurin.}

\begin{marginfigure}
\includegraphics[width=\linewidth]{fig_01}
\end{marginfigure}
Les propulseurs utilisés dans les applications militaires ou
civiles subissent, avant leur mise en service, des tests de
certification visant à contrôler leur bon fonctionnement et le
respect des normes de sécurité. Ces tests consistent à simuler
au sol les conditions de vol subies par le propulseur et à
observer les réactions de celui-ci consécutives à des
commandes de pilotage. La DGA (Direction Générale de
l'Armement) dispose dans son centre d'essais des propulseurs,
situé à Saclay, de bancs d'essais dédiés à la certification et à la
mise au point de différents types de propulseurs d'avions ou de
missiles.


Un turboréacteur est un propulseur fonctionnant sur le principe d'action-réaction. La différence de vitesse
entre l'air entrant et les gaz produits entraîne une variation de quantité de mouvement et donc un effort de
poussée (voir document 1 ci-dessous).

\begin{center}
\includegraphics[width=\linewidth]{fig_02}
\end{center}

\subsubsection*{Description structurelle du banc d’essais étudié}

Le banc d'essais de turboréacteur étudié est constitué de trois compartiments (voir document 2) :
\begin{itemize}
\item le premier compartiment (A) est alimenté par une soufflerie et a pour fonction de conditionner le flux
d'air en amont de la turbomachine testée. Il est ainsi possible de contrôler le débit, la température et
la pression de l'air en admission;
\item le deuxième compartiment (B) contient le propulseur à tester. Celui-ci est maintenu par une structure
porteuse permettant entre autres les mesures des efforts de poussée. Il est séparé du compartiment
(A) par une cloison étanche munie d'un orifice permettant le passage de l'air calibré. Le flux d'air peut
alors être laissé libre en amont du réacteur ou guidé par un raccordement jusqu'à l'entrée de celui-ci,
permettant ainsi des essais dits en "veine forcée";
\item le troisième compartiment (C) permet la collecte et l'évacuation des gaz produits lors de la
combustion.
\end{itemize}

\begin{center}
\includegraphics[width=\linewidth]{fig_03}
\end{center}

La pression à l'intérieur du compartiment B est régulée afin de simuler différentes conditions d'altitude. Des
vannes inter-compartiments permettent d'assurer une circulation d'air autour du réacteur afin de simuler le
refroidissement externe du moteur en fonctionnement. La pression du compartiment A est ajustable de 0,05 à
3 bar. Celle des compartiments B et C de 0,05 à 1,05 bar. La température d'alimentation du compartiment A
est variable de -56°C à +150°C. La capacité de ventilation est réglable de 27 à 40 kg/s. En réglant ces différents
paramètres, il est possible de simuler sur ce type de banc l'ensemble des conditions d'utilisation d'un
turboréacteur.

\question{Compléter, document réponse 1, le diagramme de définition de bloc du banc d’essais de turboréacteur.}


\subsubsection*{Description structurelle de la tuyère à ouverture variable étudiée}

Un banc d'essais nécessite pour fonctionner correctement une phase de calibration permettant d'affiner les
réglages utilisés lors des tests et d'étalonner les appareils de mesures. On s'assure notamment dans cette
phase que le compartiment A possède un comportement conforme aux besoins des tests. Les coûts en
carburant et en matériel liés à l'utilisation d'un turboréacteur sont tels que, pour ces phases de calibration, les
ingénieurs de la DGA ont imaginé une solution consistant à remplacer le propulseur réel par une structure
simulant sa présence (voir document 3).



\begin{marginfigure}
\includegraphics[width=\linewidth]{fig_04}
\end{marginfigure}

Cette structure est composée d'un tube
représentant le corps du réacteur et d'une tuyère à
ouverture variable actionnée par quatre vérins
hydrauliques et permettant de faire varier la vitesse
de l'air éjecté. On notera que dans ce cas, il n'y a pas
de combustion interne au dispositif. Le tube est fixé
sur la structure porteuse réelle avec les mêmes
points d'encrage que le propulseur et est raccordé
directement à la veine forcée. La tuyère à ouverture
variable montée sur le tube, en aval de
l'écoulement, a pour fonction de faire varier la
section de la veine de fluide en sortie de tube.

\begin{marginfigure}
\includegraphics[width=\linewidth]{fig_05}
\end{marginfigure}


La solution imaginée consiste à disposer seize volets
articulés sur la périphérie du tube qui permettent
ainsi de réduire la section de passage du fluide (voir
documents 4 et 5 ci-dessous). Ces volets sont mis en
mouvement par seize biellettes toutes identiques
reliées à une pièce de forme torique (tore) elle-même
mise en translation par quatre vérins hydrauliques
répartis régulièrement autour du tube.
Les commandes de ces vérins sont synchronisées et
asservies en position. La DGA a confié la réalisation
de cette commande à la société Bosch-Rexroth.


\begin{center}
\includegraphics[width=\linewidth]{fig_06}
\end{center}

La consigne d'ouverture de la tuyère est élaborée au niveau de la console de pilotage. Elle est transmise à des
des servo-distributeurs hydrauliques à commande électrique associés à chaque vérin. Un contrôle de la
position est effectué par un capteur à magnétostriction intégré dans le corps du vérin.

\question{Compléter le diagramme chaîne de puissance/chaîne d’information de la tuyère à ouverture variable pour banc d’essais de turboréacteur document réponse 2.}

\begin{center}
\includegraphics[width=\linewidth]{fig_07}
\end{center}


\begin{center}
\includegraphics[width=\linewidth]{fig_08}
\end{center}


